\documentclass[paper=a4,fontsize=11pt,parskip=half,subject=underlined,
               fromrule=off,fromalign=left,backaddress=false,foldmarks=false,enlargefirstpage=true]{scrlttr2}
\usepackage[T1]{fontenc}
\usepackage{lmodern}
\usepackage{microtype}
\usepackage{xcolor}
\usepackage[hidelinks]{hyperref}

\definecolor{brand}{HTML}{7A1F2B}   % the institution's colour: used for the letterhead only

% ---- Who is writing: change these once ---------------------------------------------
\newcommand*{\institution}{Example University}
\newcommand*{\department}{Department of Mathematics and Statistics}
\newcommand*{\senderline}{\textbf{Prof.\ Alex Morgan}\\ Professor of Statistics}
\newcommand*{\addressline}{10 University Road, Springfield 12345}
\newcommand*{\contactline}{+1 555 0100 \textperiodcentered\ \href{mailto:alex.morgan@example.edu}{alex.morgan@example.edu}}

% Line the address block, letterhead and signature up with the body text.
\setplength{toaddrhpos}{\dimexpr 1in+\oddsidemargin\relax}
\setplength{firstheadhpos}{\dimexpr 1in+\oddsidemargin\relax}
\setplength{firstheadwidth}{\textwidth}
\setplength{sigindent}{0pt}
\setplength{sigbeforevskip}{1.6cm}
\renewcommand*{\raggedsignature}{\raggedright}

% ---- Letterhead (first page) and footer ----------------------------------------------
% scrlttr2 draws the first page from "variables". Here the letterhead is one box:
% institution on the left, the writer on the right, and a thin rule underneath.
\setkomavar{firsthead}{%
  \parbox[b]{.52\textwidth}{\raggedright
    {\LARGE\bfseries\color{brand}\institution}\\[2pt]
    {\color{brand}\department}}%
  \hfill
  \parbox[b]{.46\textwidth}{\raggedleft\small
    \senderline\\ \addressline\\ \contactline}%
  \\[3pt]\color{brand}\rule{\textwidth}{0.8pt}}
\setkomavar{firstfoot}{\parbox{\textwidth}{\centering\footnotesize\color{black!60}
  \institution\ \textperiodcentered\ \department}}

% Later pages: a plain running head with the writer and the page number.
\pagestyle{headings}
\setkomavar{fromname}{Alex Morgan}
\setkomavar{place}{Springfield}
\setkomavar{date}{\today}
\setkomavar{subject}{Letter of recommendation for Jamie Lee}
\setkomavar{signature}{Prof.\ Alex Morgan\\ Professor of Statistics, \institution}

% Run-in heading: a bold phrase that opens a paragraph. Committees skim, so the
% phrase tells them what the paragraph is about before they read it.
\newcommand*{\topic}[1]{\textbf{#1.}\ }

\begin{document}

% The address block: who receives the letter. Use the committee's name if you do not
% know a person; "To whom it may concern" is the weakest option.
\begin{letter}{Graduate Admissions Committee\\
               Department of Computer Science\\
               Lakeshore University\\
               500 Lake Drive\\
               Lakeshore 54321}

\opening{Dear Members of the Admissions Committee,}

\topic{Summary} I am writing to recommend Jamie Lee, without reservation, for admission to
your PhD programme. In eight years of teaching and supervising I have rarely met a student
who combines such strong technical ability with such good judgement about which problems
are worth working on. State your overall verdict in the first paragraph; the rest of the
letter is the evidence for it.

\topic{How I know Jamie} I have known Jamie for three years. Jamie took my graduate course
\emph{Statistical Learning} (45 students) in autumn 2024, then wrote a master's thesis under
my supervision from February to August 2025, and has been a member of my research group
since. We meet weekly, so I can speak to Jamie's work directly, not only from grades.
Say how long you have known the candidate, in which role and how often you meet; readers
weigh the letter by how well you know the person.

\topic{Research ability} Jamie's thesis extended a recent method for graph estimation to
data with heavy-tailed noise, a problem I had set aside as too hard. Jamie found a
simpler proof route than the one I had in mind, implemented the method independently, and
showed that it beats the best published baseline on three of four benchmarks. The work is
now a conference submission with Jamie as first author. Concrete examples like this one,
with what the candidate did and what came of it, persuade more than adjectives.

\topic{Independence and communication} When a collaborator's data turned out to be
mislabelled, Jamie noticed it, traced the cause and proposed a fix before anyone
asked. Jamie also explains difficult ideas clearly: the group seminar talk was the best
given by a master's student in my time here, and the students in the tutorial Jamie
co-taught rated it 4.8 out of 5.

\topic{Comparison} Of the roughly 90 students I have supervised or taught closely, Jamie
is among the top three, and the strongest in the last five years. Compared with
applicants I have recommended to your programme before, Jamie is on a par with the best.
Give a ranking or a percentage and name the group you compare with; ``excellent'' alone
says little, because every letter says it.

\topic{Limits} In fairness, Jamie is quiet in large meetings and sometimes needs
encouragement to share ideas early. In small groups this is no obstacle. A short, honest
remark like this makes the rest of the letter more believable.

\topic{Recommendation} I recommend Jamie Lee in the strongest terms and am confident that
Jamie will do excellent research in your programme. Please contact me at
\href{mailto:alex.morgan@example.edu}{alex.morgan@example.edu} if I can add anything.

\closing{Sincerely,}

% Optional: what you attach, or who gets a copy.
% \encl{Curriculum vitae of the candidate}

\end{letter}
\end{document}
